\documentclass[11pt,a4paper]{article}

\usepackage[spanish]{babel}
\usepackage[utf8]{inputenc}
\usepackage[T1]{fontenc}
\usepackage{lmodern}
\usepackage{amssymb}
\usepackage[a4paper,margin=2.4cm]{geometry}
\usepackage{xcolor}
\usepackage{titlesec}
\usepackage{enumitem}
\usepackage{booktabs}
\usepackage{fancyhdr}
\usepackage[hidelinks]{hyperref}

% ---- Paleta ----
\definecolor{indigo}{HTML}{312E81}
\definecolor{indigol}{HTML}{4338CA}
\definecolor{teal}{HTML}{0E7C7B}
\definecolor{coral}{HTML}{C2410C}
\definecolor{ink}{HTML}{1F2430}
\definecolor{rule}{HTML}{C9C6BE}

% ---- Títulos ----
\titleformat{\section}{\large\bfseries\color{indigo}}{\thesection.}{0.6em}{}
\titlespacing*{\section}{0pt}{1.4em}{0.5em}
\titleformat{\subsection}[runin]{\bfseries\color{indigol}}{}{0em}{}[\hspace{0.4em}]
\titlespacing*{\subsection}{0pt}{0.9em}{0.4em}

\setlist{nosep,leftmargin=1.5em}
\setlength{\parindent}{0pt}
\setlength{\parskip}{0.35em}

% ---- Bloque de "guion hablado" ----
\newenvironment{guion}%
{\par\smallskip\begingroup\leftskip=1.4em\rightskip=0.8em\color{ink}\itshape}%
{\par\endgroup\smallskip}

\newcommand{\lab}[1]{\textbf{\color{ink}#1}\ }

% ---- Encabezado/pie ----
\pagestyle{fancy}\fancyhf{}
\lhead{\small\color{ink}Sueño, memoria y señales fisiológicas}
\rhead{\small\color{ink}Guion de exposición}
\cfoot{\small\thepage}
\renewcommand{\headrulewidth}{0.4pt}

\begin{document}

% ================= PORTADA =================
\begin{center}
  {\LARGE\bfseries\color{indigo} Guion de exposición}\\[0.35em]
  {\large Sueño, memoria y señales fisiológicas}\\[0.7em]
  {\normalsize Trabajo Práctico \textemdash{} Neurociencias del Desarrollo Productivo}\\[0.3em]
  {\small Keoni Dubovitsky \quad\textbullet\quad Joaquín Pelufo \quad\textbullet\quad Sergio Tenoch Sil Cruz}\\[0.2em]
  {\small Fecha de exposición: 18/6/2026 \quad\textbullet\quad Duración objetivo: 12--15 min}
\end{center}
\vspace{0.4em}
{\color{rule}\hrule}
\vspace{0.8em}

\noindent\small
El presente documento reúne el guion de la exposición oral del trabajo. El texto marcado como
\textit{Guion} está redactado en registro académico oral (primera persona del plural, terminología
formal), de modo que la exposición conserve el tono de una comunicación de tipo \textit{paper}.
Para la presentación en vivo, abrir \texttt{index.html} y presionar \texttt{S} (vista de orador):
muestra estas notas, el cronómetro y la diapositiva siguiente.
\normalsize

% ================= REPARTO =================
\section{Reparto de oradores}

\begin{center}
\begin{tabular}{@{}lll@{}}
\toprule
\textbf{Bloque} & \textbf{Orador} & \textbf{Diapositivas} \\
\midrule
Apertura, encuadre y marco & Keoni Dubovitsky & 1--8 \\
Métodos, memoria y calidad de señal & Joaquín Pelufo & 9--16 \\
Ondas, contexto y cierre & Sergio Tenoch Sil Cruz & 17--24 \\
\bottomrule
\end{tabular}
\end{center}

\noindent\textit{Transiciones:} al concluir el último slide del bloque, ceder la palabra nombrando al
siguiente orador. Una idea por diapositiva; no leer la pantalla. Convención: \lab{Idea.} la noción que
debe quedar establecida; \lab{Guion.} texto sugerido para la exposición; \lab{Cuidado.} precisión
metodológica o terminológica a preservar.

% ================= BLOQUE 1 =================
\section{Bloque 1 — Keoni (diapositivas 1--8)}

\subsection{Slide 1 · Portada.}
\lab{Idea.} Identificación del equipo y tesis del trabajo.
\begin{guion}
«Buenos días. Presentamos nuestro trabajo sobre sueño, memoria y señales fisiológicas. La pregunta que
lo guía es si el sueño posterior al aprendizaje favorece la consolidación de la memoria declarativa.
Adelantamos la tesis: la evidencia obtenida resulta \emph{compatible} con dicha hipótesis, en carácter
de apoyo descriptivo y parcial, sin constituir una demostración causal.»
\end{guion}
\lab{Cuidado.} Establecer desde el inicio el registro «compatible / parcial».

\subsection{Slide 2 · Síntesis del estudio.}
\lab{Idea.} Resultado principal y su límite, de entrada.
\begin{guion}
«En síntesis: el sujeto de la condición de sueño obtuvo 17 sobre 20 en la evaluación final, y el
registro de EEG alcanzó sueño NREM profundo hasta S4. No obstante, la condición de sueño cuenta con un
único sujeto y el EEG proviene de un registro independiente; por ello la relación se plantea como
integración teórica, no como correlación directa.»
\end{guion}

\subsection{Slide 3 · Objetivo e hipótesis.}
\lab{Idea.} Objetivo, hipótesis y criterio fijado a priori.
\begin{guion}
«El objetivo es evaluar si el sueño favorece la consolidación de la memoria declarativa, integrando el
desempeño conductual y la fisiología del sueño. La hipótesis postula que el sueño mejora dicha
consolidación. El criterio de decisión se fijó a priori: la hipótesis se considera apoyada únicamente
si el patrón conductual y la arquitectura fisiológica resultan compatibles.»
\end{guion}
\lab{Cuidado.} Enfatizar «criterio definido a priori».

\subsection{Slide 4 · Predicciones.}
\lab{Idea.} Predicciones y resultado esperado bajo sueño (lo solicita la consigna).
\begin{guion}
«Se definieron tres predicciones a priori: mejor recuerdo en la condición de sueño, una arquitectura
NREM con sueño profundo y la presencia de los ritmos de consolidación. Respecto de la consigna: ante una
siesta consolidada de unos noventa minutos cabría esperar al menos un ciclo con sueño profundo y,
posiblemente, una ventana de REM. El sujeto propio no concilió el sueño, circunstancia que se retoma como
limitación.»
\end{guion}

\subsection{Slide 5 · Diseño.}
\lab{Idea.} Tres etapas, dos condiciones; $n=1$ en sueño.
\begin{guion}
«El diseño consta de tres etapas: codificación de veinte asociaciones palabra--definición (TR); un
intervalo en que una condición duerme una siesta de unos noventa minutos y la otra permanece en vigilia;
y una evaluación final (TS). Se analizan el desempeño final y el cambio TR--TS. Cabe señalar que la
condición de sueño cuenta con un único sujeto, por lo que la lectura es descriptiva.»
\end{guion}
\lab{Cuidado.} Explicitar $n=1$ antes de que lo señalen.

\subsection{Slide 6 · Tres fuentes de datos independientes.}
\lab{Idea.} La regla metodológica central.
\begin{guion}
«Los datos provienen de tres fuentes independientes: la tarea de memoria, de una planilla propia con
cuatro sujetos; el EEG de sueño, un registro secundario de la cátedra correspondiente a una persona que
sí durmió; y un registro propio con OpenBCI, cuyo sujeto no concilió el sueño. La regla metodológica
central es que memoria y EEG corresponden a sujetos distintos y no se correlacionan de forma directa.»
\end{guion}
\lab{Cuidado.} Ante preguntas, reiterar: sujetos distintos.

\subsection{Slide 7 · Procedimiento (imágenes).}
\lab{Idea.} Documentar el procedimiento (requisito de la consigna).
\begin{guion}
«El procedimiento, conforme lo solicita la consigna. Se colocó un montaje central C3--C4 con
electrooculograma y electromiograma —una polisomnografía reducida—; se preparó el entorno y el sujeto
quedó dispuesto para la siesta. Este registro propio es el que posteriormente no logró conciliar el
sueño.»
\end{guion}

\subsection{Slide 8 · Marco teórico.}
\lab{Idea.} Consolidación activa $=$ ondas lentas $+$ husos $+$ ripples.
\begin{guion}
«El marco teórico es la consolidación activa de sistemas. Durante el sueño profundo se acoplan tres
ritmos —ondas lentas, husos y ripples hipocampales—, lo que reactiva memorias dependientes del hipocampo
y favorece su integración cortical. La sigma se emplea como \emph{proxy} de husos. Asimismo, lo
determinante no es la mera profundidad del sueño, sino el acople fino entre estos ritmos.»
\end{guion}
\lab{Cuidado.} Reiterar «proxy» en cada mención de sigma.

% ================= BLOQUE 2 =================
\section{Bloque 2 — Joaquín (diapositivas 9--16)}

\subsection{Slide 9 · Métodos.}
\lab{Idea.} Dos análisis en planos independientes.
\begin{guion}
«Se aplicaron dos análisis. Para memoria, dos puntajes: la coincidencia formal exacta de la pseudopalabra
y el contenido semántico de la definición. Para el EEG, se aplicaron filtros previos —notch de 50 Hz y
pasa-banda de 0,3 a 35 Hz— y se segmentó en épocas de treinta segundos. El canal C4 presentaba artefactos,
de modo que el análisis se realizó sobre C3.»
\end{guion}

\subsection{Slide 10 · Resultado conductual (sueño vs.\ control).}
\lab{Idea.} Mayor desempeño en sueño, sobre una base elevada.
\begin{guion}
«El sujeto de la condición de sueño alcanzó el mayor desempeño final, 17 sobre 20. No obstante, su línea
de base era elevada —15—, por lo que el incremento neto fue de dos puntos. El grupo control en vigilia
promedió un desempeño bastante menor (3,7 a 5,0), aunque uno de sus sujetos presentó el mayor incremento
individual, de 10 a 14. El resultado es compatible con la hipótesis, pero con un único sujeto no permite
inferir causalidad.»
\end{guion}
\lab{Cuidado.} Mencionar proactivamente que un sujeto en vigilia tuvo el mayor incremento.

\subsection{Slide 11 · Detalle conductual por sujeto.}
\lab{Idea.} Las dos métricas y por qué la semántica no discrimina.
\begin{guion}
«En detalle: la métrica de definición —semántica— está próxima al techo en todos los sujetos, por lo que
no discrimina entre condiciones. Toda la señal recae en la palabra objetivo, la coincidencia formal
exacta, que es más exigente. Con cuatro sujetos y alta heterogeneidad, lo leemos como descriptivo.»
\end{guion}

\subsection{Slide 12 · EEG: control de calidad.}
\lab{Idea.} Preprocesamiento previo a la interpretación; descarte de C4.
\begin{guion}
«El preprocesamiento precede a la interpretación. El canal C3 presentó una calidad adecuada; C4 mostró
una desviación dieciocho veces mayor y saturación, de modo que su promedio contaminaba el espectro. En
consecuencia, se reporta únicamente C3.»
\end{guion}

\subsection{Slide 13 · Estadificación.}
\lab{Idea.} Clave de la cátedra; el registro alcanzó S4.
\begin{guion}
«La estadificación se encuentra en la primera columna del archivo, una fase por época. Conforme a la clave
de la cátedra, el código 2 corresponde a S3 y el 3 a S4, la fase más profunda. No figuran los códigos 4, 5
ni 8, de modo que no se registró REM. El conjunto S3 y S4 representa cerca del 69\,\% del registro: se
alcanzó sueño profundo.»
\end{guion}
\lab{Cuidado.} Esta clave no codifica un S2 independiente; explicitarlo, y por ello no afirmar la detección de husos de N2.

\subsection{Slide 14 · Hipnograma.}
\lab{Idea.} Descenso NREM sostenido hasta S4; ausencia de REM.
\begin{guion}
«El hipnograma describe la evolución temporal. La latencia de sueño fue de cuatro minutos; el registro
alcanzó S4 a los 45,5 minutos y sostuvo ese bloque profundo hasta el minuto 68. No se observó REM.
Predomina el NREM, con una notable profundidad.»
\end{guion}

\subsection{Slide 15 · Distribución de fases.}
\lab{Idea.} Cerca del 28\,\% en la fase más profunda; eficiencia del 86\,\%.
\begin{guion}
«En cuanto a la distribución, S3 fue la fase predominante, con un 41\,\%, y S4 aportó un 28\,\% adicional;
en conjunto, cerca del 69\,\% correspondió a sueño profundo. El sueño total fue de 68 minutos, con una
eficiencia del 86\,\%: un registro de buena calidad.»
\end{guion}

\subsection{Slide 16 · Firma espectral.}
\lab{Idea.} La delta aumenta con la profundidad; la sigma se concentra en NREM (proxy).
\begin{guion}
«Respecto de la firma espectral, la potencia delta aumenta con la profundidad y alcanza el 91,6\,\% en S4,
en concordancia con la teoría. La sigma, proxy de husos, presenta su máximo en S3. Conviene precisar que
los husos son, teóricamente, propios de N2; dado que esta clave no codifica un N2 independiente, se reporta
sigma elevada en NREM y no la detección formal de husos.»
\end{guion}
\lab{Cuidado.} «Potencia de banda, no detección formal de husos».

% ================= BLOQUE 3 =================
\section{Bloque 3 — Sergio (diapositivas 17--24)}

\subsection{Slide 17 · Morfología del EEG por fase (animación).}
\lab{Idea.} Representar la firma espectral como forma de onda.
\begin{guion}
«Para ilustrar la firma espectral, estas cuatro trazas se sintetizaron a partir de la potencia de banda
real de cada fase. En vigilia, actividad rápida y de baja amplitud; en S1, amplitud reducida; en S3, ondas
lentas con husos superpuestos —donde la sigma fue máxima—; y en S4, ondas lentas de gran amplitud que
dominan el trazado.»
\end{guion}
\lab{Cuidado.} Aclarar que son trazas reconstruidas, no el registro crudo.

\subsection{Slide 18 · Espectrograma.}
\lab{Idea.} Confirmación visual de la firma espectral sobre el eje temporal.
\begin{guion}
«El espectrograma muestra la potencia por frecuencia a lo largo del registro, con el hipnograma alineado
en la parte superior. La potencia lenta se concentra en el bloque de sueño profundo, entre los minutos 45
y 68, mientras la actividad rápida se atenúa. Es la misma firma espectral, ahora sobre el eje temporal.»
\end{guion}
\lab{Cuidado.} Representación descriptiva sobre C3; la sigma, proxy de husos.

\subsection{Slide 19 · Análisis espectral — detalle.}
\lab{Idea.} PSD por fase y composición de bandas (respaldo cuantitativo).
\begin{guion}
«Dos análisis complementarios. La densidad espectral de potencia por fase muestra, en S3, un realce en la
banda sigma y, en S4, el predominio de las frecuencias lentas. La composición relativa de bandas confirma
que el predominio de delta aumenta progresivamente con la profundidad.»
\end{guion}

\subsection{Slide 20 · Ficha técnica.}
\lab{Idea.} Parámetros y terminología para sostener la metodología.
\begin{guion}
«A modo de ficha técnica: registro BrainVision de diez canales a 250 Hz; canal C3, con C4 descartado;
filtros notch de 50 Hz y pasa-banda de 0,3 a 35 Hz; 159 épocas de 30 segundos. Las bandas y fases
observadas se resumen aquí, junto con el lenguaje que adoptamos: compatible, apoyo descriptivo, proxy.»
\end{guion}
\lab{Cuidado.} Mostrarla con agilidad; es soporte, no se lee entera.

\subsection{Slide 21 · Alcance de las conclusiones.}
\lab{Idea.} Delimitar lo sostenible y lo no sostenible.
\begin{guion}
«En cuanto al alcance: los datos permiten afirmar la presencia de NREM profundo, una potencia delta máxima
en S4 y sigma elevada, compatibles con los ritmos de consolidación. No permiten afirmar que dichos husos
expliquen la memoria de nuestros sujetos —se trata de personas distintas— ni que exista una relación
causal. La relación memoria--EEG se plantea como integración teórica.»
\end{guion}

\subsection{Slide 22 · Factores que pudieron impedir el sueño.}
\lab{Idea.} Moduladores del descanso (integra varios contenidos del curso).
\begin{guion}
«Respecto de los factores que pudieron impedir la conciliación del sueño: una siesta vespertina puede
situarse en contrafase del cronotipo, y la exposición lumínica retrasa la fase circadiana; la cafeína
antagoniza la adenosina y eleva la latencia; el estrés y el cortisol dificultan la conciliación; y el
registro en laboratorio, con electrodos, incrementa la incomodidad. Son variables plausibles, no causas
medidas.»
\end{guion}

\subsection{Slide 23 · Estrés, eje HPA y consolidación.}
\lab{Idea.} Mecanismo por el cual el estrés previo al sueño degrada la consolidación (consigna).
\begin{guion}
«La consigna pregunta por el efecto de un evento estresante antes de dormir. El estrés activa el eje
hipotálamo--hipófisis--suprarrenal: el hipotálamo libera CRH, la hipófisis ACTH y la suprarrenal cortisol.
Lo determinante para la memoria es que el cortisol elevado durante el sueño profundo temprano perjudica la
consolidación declarativa dependiente del hipocampo (Plihal \& Born). En nuestro caso, un evento estresante
previo a la siesta dificultaría conciliar el sueño y degradaría el sueño profundo y los husos que sostienen
la consolidación.»
\end{guion}
\lab{Cuidado.} Mantener el crédito de la imagen (Shana Cannella, CC BY 3.0).

\subsection{Slide 24 · Conclusión.}
\lab{Idea.} Respuesta directa y requisitos para una afirmación más sólida.
\begin{guion}
«En conclusión, los datos son compatibles con la hipótesis, sin constituir una demostración causal. Como
apoyo descriptivo: el mayor desempeño en la condición de sueño y un EEG con sueño profundo. Para una
afirmación más sólida se requeriría medir memoria y EEG en un mismo sujeto, un mayor tamaño muestral, el
control de los moduladores y un detector formal de husos. El sueño se asocia a la consolidación de la
memoria declarativa, con las reservas que los datos imponen. Muchas gracias.»
\end{guion}

% ================= Q&A =================
\section{Preguntas anticipadas}

\begin{enumerate}[label=\textbf{\arabic*.},leftmargin=2em]
\item \textbf{¿El sujeto del EEG es uno de los de la tarea de memoria?}\\
Corresponden a fuentes y sujetos distintos: la memoria proviene de la planilla propia; el EEG, de un
registro secundario de la cátedra de una persona que durmió. Por ello no se computa correlación
memoria--EEG: la relación es de integración teórica.

\item \textbf{Con $n=1$ en la condición de sueño, ¿la evidencia no es débil?}\\
En efecto, y se asume explícitamente: por eso se habla de un patrón descriptivo «compatible», no de una
prueba. El aporte reside en integrar conducta, fisiología y teoría de manera rigurosa.

\item \textbf{¿Por qué la clave indica $2=$S3 y $3=$S4, si lo habitual es $2=$S2, $3=$S3?}\\
Se emplea la clave confirmada por la cátedra, que no codifica un S2 independiente y rotula $2=$S3 y $3=$S4.
La conclusión se sostiene igualmente: se alcanzó sueño profundo y no se registró REM (código 5 ausente).

\item \textbf{¿Husos en S3? ¿No deberían situarse en N2?}\\
Teóricamente, los husos son propios de N2. Dado que esta clave no contempla un N2 independiente, la sigma
se reporta como potencia de banda elevada en NREM —un proxy—, no como husos detectados. Un detector formal
(p.\ ej.\ YASA) constituye trabajo futuro.

\item \textbf{¿Por qué se descartó C4?}\\
C4 presentó una desviación dieciocho veces superior a la de C3 y saturación; su promedio contaminaba el
espectro. Reportar únicamente C3 es la decisión conservadora.

\item \textbf{¿Por qué no se observa REM, si la siesta fue de unos 90 minutos?}\\
Cabía esperar REM, pero el registro útil (79,5 minutos) permaneció en NREM profundo; el código de REM no
figura. Es esperable, dado que el REM suele aparecer en ciclos posteriores.

\item \textbf{¿Cómo incidirían la cafeína o el estrés previos al sueño? (consigna)}\\
La cafeína antagoniza la adenosina, eleva la latencia y fragmenta el sueño, reduciendo SWS y husos. El
estrés, mediante cortisol elevado en el SWS temprano, perjudica la consolidación hipocampal (Plihal \&
Born), además de dificultar la conciliación.

\item \textbf{¿De dónde proceden las ondas animadas?}\\
Son trazas sintetizadas a partir de la potencia de banda real medida en C3 por fase (delta, theta, alfa,
sigma, beta). No corresponden a la señal cruda: ilustran la firma espectral en forma de onda.
\end{enumerate}

% ================= CIERRE =================
\section{Cierre metodológico — criterios no negociables}

\begin{itemize}[leftmargin=1.5em]
\item Emplear: \textbf{«compatible con»}, \textbf{«apoyo descriptivo / parcial»}, \textbf{«proxy de sigma»}.
\item Evitar: \textbf{«demostrado»}, \textbf{«causal»}, \textbf{«los husos explican el resultado»}.
\item Memoria y EEG: \textbf{fuentes distintas}; no correlación directa.
\item Condición de sueño: $\mathbf{n=1}$; explicitarlo proactivamente.
\item Sigma: potencia de banda, \textbf{no} detección formal de husos.
\end{itemize}

% ================= CHECKLIST =================
\section{Checklist de ensayo}

\begin{itemize}[leftmargin=1.5em]
\item[$\square$] Ensayo cronometrado completo $\leq 15$ min (cada bloque $\leq 5$ min).
\item[$\square$] Cada orador practica sus transiciones y cede la palabra al siguiente.
\item[$\square$] Probar la vista de orador (\texttt{S}) en el equipo de la exposición.
\item[$\square$] Verificar el funcionamiento sin conexión y a pantalla completa (\texttt{F}).
\item[$\square$] Confirmar que las animaciones de onda se reproducen con fluidez en el proyector.
\item[$\square$] Disponer de un PDF de respaldo ante una eventual falla del navegador.
\item[$\square$] Repasar las ocho preguntas anticipadas y asignar el responsable de cada tema.
\end{itemize}

\end{document}
